\section{结论与局限}

本文研究少量双臂示范下的协同时序数据生成问题，并提出\method{}。该方法通过动作完整的单臂后缀延迟构造受控反事实，在物理仿真中识别双臂协同敏感阶段，再将成功近邻和困难失败样本用于策略学习。其核心出发点是：运动学上可达并不保证物理上可执行，而相对时序是接触丰富双臂任务中不可忽略的变量。

当前方案仍有四点局限。第一，插入等待虽然保留了动作序列，但系统状态在等待期间发生变化，因此它是面向执行因果性的干预，而不是严格保持所有其他变量不变的统计反事实。第二，delta 动作完整并不等价于末端状态完全一致，控制器闭环误差和接触动力学仍会改变最终状态。第三，仿真接触模型和真实机器人之间存在差距，仿真中得到的敏感边界需要少量实机试验校准。第四，单一任务上的非单调响应可能来自特定接触几何，必须通过更多任务和初始条件验证其普遍性。

后续工作将研究从第一视角人类视频或跨具身重定向轨迹中恢复阶段语义，并用语义阶段约束替代固定比例窗口；同时探索使用视觉语言模型辅助识别“支撑、对齐、插入、释放”等协同阶段。\todo{根据最终实验结果重写本节，避免重复摘要。}
